\section{Proposed Methodology}

\subsection{System View}

The proposed system is a research prototype and benchmark. It is not presented
as a production deployment. The design has four main parts: an access request
gateway, a layered access-control service, an explainable decision service, and
a permissioned blockchain-style audit layer. Raw records remain in
agency-controlled off-chain storage. The audit layer stores only identifiers,
commitments, hashes, policy versions, approval references, explanation hashes,
and block metadata.

\subsection{Access Request Flow}

The workflow begins when a user requests a sensitive record. The user may be a
police officer, an investigating agency user, a forensic user, a prosecutor or
court-related authority, or an auditor. The gateway converts the request into a
structured form. Subject attributes include role, rank, station, jurisdiction,
credential status, and case assignment. Object attributes include record type,
sensitivity level, juvenile or witness flags, sealed status, and owning agency.
Action attributes describe what the requester wants to do. Environment
attributes include purpose, time window, emergency flag, court/prosecutor link,
approval token, and policy version.

The access-control service evaluates these attributes and returns
\emph{allow}, \emph{deny}, or \emph{escalate}. A request can be allowed when
the requester has the required role, assignment, purpose, and approval context.
It can be denied when the credential is revoked, the approval is invalid, or
the request clearly violates policy. It can be escalated when a superior or
review authority must check the request, such as in cross-jurisdiction or
sensitive-record cases.

\subsection{Policy Layer}

The policy layer combines RBAC, ABAC, and PBAC-style rules. RBAC gives the
basic user categories. ABAC adds contextual attributes such as jurisdiction,
case assignment, purpose, sensitivity, and credential state. PBAC adds policy
versioning, approval requirements, revocation handling, and exceptional cases.
The current prototype uses a deterministic policy oracle because the first
paper needs transparent benchmark labels. This is more suitable than a black
box model for the initial access decision.

\subsection{Explainable Decision Layer}

After a decision is made, the XAI layer creates an explanation artifact. The
artifact contains the decision label, reason code, decisive attributes, policy
or rule version, counterfactual information where applicable, and an
explanation hash. A typical explanation answers a direct question: why was this
request allowed, denied, or escalated?

This layer is designed for review, not only for display. A supervisor can use
it to inspect an escalation. An auditor can use it to check whether the policy
and explanation match. A court-related or administrative reviewer can use it to
understand what information was considered at the time of access. The
explanation hash is linked to the audit event so that a later change in the
explanation can be detected.

\subsection{Blockchain-Style Audit Layer}

The audit layer is modeled as a permissioned blockchain-style ledger among
known agencies. Each audit event stores minimized metadata and commitments
rather than raw sensitive records. A block contains a batch of events, a
previous-block hash, an event commitment, validator identity, and signature.
The design is compared with simpler baselines such as a mutable log and a
signed append-only hash chain. This comparison is important because blockchain
should not be assumed to help unless it is tested against simpler audit
mechanisms.

\subsection{Synthetic Benchmark}

Real police access logs are not publicly available, and they should not be
assumed in a first-stage academic prototype. The benchmark therefore generates
synthetic stations, officers, cases, records, and access requests. The declared
policy oracle labels each request as allow, deny, or escalate. The synthetic
setting gives control over request context, policy stress cases, and attack
conditions. It does not claim to represent real CCTNS or ICJS traffic.

\subsection{Threat Model and Evaluation}

The attack catalog includes ordinary tampering, approval-token replay, request
backdating, explanation-hash substitution, revocation-race behavior,
metadata-inference checks, block-signature collusion, and the validly re-signed
compromised-signer attack. The compromised-signer case is important because an
attacker may corrupt the policy output and then produce a log that still looks
properly signed. This separates integrity of the log from correctness of the
decision.

The evaluation compares mutable logs, signed hash chains, blockchain-style
audit, CT-style logs, ABAC/Fabric-style re-execution, a trusted raw-attribute
policy oracle, and NS-PI log-only drift detection. The research questions are:

\begin{itemize}
    \item RQ1: Which audit designs detect ordinary tamper attacks?
    \item RQ2: Which defenses detect validly re-signed policy corruption?
    \item RQ3: What are the sensitivity limits of NS-PI drift detection?
    \item RQ4: How complete and stable are the XAI artifacts and audit
    reconstruction traces?
    \item RQ5: What metadata exposure and local latency/storage overhead are
    introduced by the audit design?
\end{itemize}

The metrics include adversarial audit score, per-attack detection rate,
trace-completeness rate, decisive-attribute coverage, counterfactual coverage,
counterfactual validity, explanation stability, audit reconstruction rate,
metadata exposure score, decision/build latency, verification latency, and
storage overhead. The experiments are run with deterministic seeds
\{7, 21, 42, 99, 123\}. The paper reports results only from generated
artifacts in the repository and does not infer real police-system performance
from the synthetic workload.
